\documentclass[12pt, a4paper]{article}

% --- Packages ---
\usepackage[margin=1.8cm]{geometry}
\usepackage{amsmath}
\usepackage[hidelinks]{hyperref}
\usepackage[T1]{fontenc}
\usepackage{lmodern}
\usepackage{microtype}
\usepackage{setspace}
\usepackage{parskip}
\usepackage{natbib}

\setlength{\parskip}{0pt}
\setlength{\parindent}{1.5em}

\begin{document}

% -------------------------------------------------------
%  TITLE BLOCK
% -------------------------------------------------------
\begin{center}
    {\LARGE The Political Legacy of Displacement:\\[4pt]
    Evidence from the Spanish Exile}\\[18pt]
    {\large Cristina Aranzana$^{*}$, and Luc Paluskiewicz$^{\dagger}$}\\[10pt]
\end{center}

\vspace{10pt}

% -------------------------------------------------------
%  ABSTRACT
% -------------------------------------------------------
\begin{abstract}
\setstretch{1.1}

As of the end of June 2025, 117.3 million people had been forced to flee their homes
globally due to violence, conflict, and human rights violations (UNHCR, 2025).
When persecution is politically motivated, exile becomes politically selective. What
are the political consequences for host communities when they receive ideologically
motivated refugees? Contact between the displaced and the native population could not
only reduce prejudice (Allport, 1954), but also facilitate the diffusion
of ideas between both groups. On the contrary, if the integration of the displaced
population is limited, the ideas refugees bring may be perceived as a specific threat,
generating political backlash. Theoretically, the expected effect is
therefore ambiguous. To examine this question, we focus on the mass exodus from Spain to France during the
final phase of the Spanish Civil War, known as \textit{La Retirada} (``the retreat'').
In early 1939, the collapse of Republican forces in Catalonia---who represented various
left-wing groups including anarchists, communists, and other republican
factions---triggered the flight of nearly half a million people. French authorities,
unprepared for the scale of the inflow, treated the refugees as a security concern and
responded by interning them. Because no permanent facilities had been prepared in
advance, camp locations were chosen under extreme time pressure based on logistical
availability. This process generated a
geographically uneven yet plausibly exogenous distribution of camp locations.

\medskip


By leveraging this historical episode, we estimate the effect of being exposed to a camp on voting behaviour using a difference-in-differences
strategy. Exposure to a refugee
camp decreased the share of votes for the socialist party by 2.7 percentage
points---a 10\% reduction relative to the pre-treatment mean. Given that the refugees
were themselves strongly aligned with left-wing and socialist ideals, this decline is
consistent with a political backlash among host populations. Yet the response was not
uniform across the political spectrum. Several classic left-wing organizations---
particularly labour and professional unions---were actively involved in assisting the
refugees and thus maintained close contact with them. Reflecting this heterogeneity,
we find a positive impact on support for the Communist Party: municipalities within
1~km of a refugee camp increased their PCF vote share by 1.6 percentage points---
a 15\% increase over the pre-treatment mean. Aligned with these results, municipalities that hosted refugees also
exhibit higher levels of resistance against the Nazi occupation.  We also observe stronger effects among younger cohorts, consistent with the idea that younger populations are more influenced by their proximity to foreign political refugees. Further heterogeneity analysis suggests that the political effects are attributable to municipalities in proximity to high-capacity internment camps. Taken together, these findings show that the arrival of politically motivated refugees
can generate political polarisation. 



\end{abstract}

\vspace{12pt}

\noindent\textbf{Keywords:} forced displacement, political backlash, idea diffusion,
refugee camps, political polarisation, difference-in-differences, spatial spillovers,
Spanish Civil War, France.

\noindent\textbf{JEL Classification:} D72, F22, J15, N44, P16.

\vfill

% --- Affiliations ---
\noindent\rule{0.4\linewidth}{0.4pt}

\noindent{\small $^{*}$University of Barcelona;\, e-mail: c.aranzana@ub.edu}

\noindent{\small $^{\dagger}$Paris School of Economics, Coll\`ege de France;\,
e-mail: luc.paluskiewicz@psemail.eu}


\end{document}